\scanpage{204}
\mspage{120}
% archive sheet number: 203
% element: text | top, line 1 | 2 Considérons
$2^\circ$) Considérons, pour $x\in\mathrm{Ob}\,C$, le foncteur
% element: diagram | top, line 2 | e to C labelled epsilon x
\diagram{\begin{tikzcd}[ampersand replacement=\&]
\unsure{e|\mathbf e}\arrow[r,"\unsure{\varepsilon_x|\epsilon_x}"]\&C,
\end{tikzcd}}
% element: text | upper third, line 3 | d'où un foncteur
d'où un foncteur
% element: display | upper third, line 4 | epsilon x star evaluation
\[
\varepsilon_{x\deleted{\unsure{*}}}^*:\underline{\mathrm{Hom}}(C,\unsure{\mathscr E|\mathcal E|E})\longrightarrow\mathscr E,\qquad F\longmapsto F(x),
\]
% element: text | upper third, line 5 | admettant un adjoint
admettant un adjoint à droite
% element: display | upper third, line 6 | epsilon x lower star
\[
\varepsilon_{x*}:\mathscr E\longrightarrow\underline{\mathrm{Hom}}(C,\mathscr E),
\]
% element: text | upper middle, line 7 | donné par
donné par
% element: display | upper middle, line 8 | epsilon x star E of y
\[
\varepsilon_{x*}(E)(\unsure{y|x})=E^{\mathrm{Hom}(x,y)}
\]
% element: text | middle, lines 9-12 | sous réserve d'existence
(sous réserve d'existence des produits de cardinal $\leq\pi$). Le foncteur $\varepsilon_{x*}$ est $\pi$-accessible en vertu de loc. cit. 9.8.

% element: text | middle, lines 13-17 | Ceci posé
Ceci posé, soit $F$ un objet $\pi$-accessible de $\underline{\mathrm{Hom}}(C,\mathscr E)$, on voit que $\varepsilon_x^*(F)$ est $\pi$-accessible, $\varepsilon_{x\deleted{\unsure{*}}}^*$ \deleted{\unsure{adjoint}} ayant un adjoint à droite $\pi$-accessible :
% element: deletion | above next display, left | Hom epsilon x
\deleted{$\unsure{\mathrm{Hom}_{\varepsilon_x}|\mathrm{Hom}(\varepsilon_x)}$}
% element: display | lower middle, line 18 | Hom adjunction formula
\[
E\longmapsto\mathrm{Hom}_{\mathscr E}\bigl(\underset{\varepsilon_x^*(F)}{\underline{F(x)}},E\bigr)
\simeq\mathrm{Hom}_{\deleted{\unsure{\mathscr E}}\underline{\mathrm{Hom}}(C,\mathscr E)}(F,\varepsilon_{x*}(F))
\]
% element: text | lower third, lines 19-22 | est composé des foncteurs
est composé \unsure{du foncteur|des foncteurs} $\pi$-\unsure{accessible|accessibles} $\varepsilon_{x*}$, et du foncteur \inserted{cov.} ($\pi$-accessible) représenté par $F$ dans $\underline{\mathrm{Hom}}(C,\mathscr E)$.

% element: text | foot, lines 23-26 | Inversement
Inversement si les $F(x)$ sont $\pi$-accessibles, \unsure{prouvons|prouver} que $F$ l'est, \inserted{moyennant $\operatorname{card}\mathrm{Fl}(C)\leq\pi$} (\unsure{se} fait par AQT), je ne vais pas l'expliciter ici.
